\paragraph{Two padding orbits and a stronger exponent constant.}
Allowing much larger splitting primes gives a different tradeoff. One
can retain unique, efficiently recoverable nearby witnesses and almost
full-gap separation even when the numerical prescription's exponent
constant is doubled. This construction does not retain logarithmic
length in the field size.

\begin{lemma}[Rigidity for nearby radical radii]\label{tou:rigidity}
Let $d$ be prime, let $q$ be a power of two with
$q-1>3d^{d-1}$, and let $M\ge1$. Choose an integer $L\ge1$ with
$q^L>4(dM)^{d-2}$. Write
\[
 t=q^{1/d}>0,\qquad
 R_i=\frac{q^{i+1}-1}{q^i-1},\qquad
 a_i=R_i^{1/d}>0\quad(L\le i<L+M),
\]
and let $\zeta$ be a complex primitive $d$th root of unity.
Every zero-sum subset of the orbits $a_i\{1,\zeta,\ldots,\zeta^{d-1}\}$
is a union of whole orbits.
The same holds if some orbits are replaced by single points: no such
single point can occur in a zero-sum subset.
\end{lemma}
\begin{proof}
Write $a_i=t+\delta_i$. Since
$R_i=q+(q-1)/(q^i-1)$, the difference-of-powers identity gives
\[
 \frac{t}{3d}q^{-i}<\delta_i\le\frac{t}{d}q^{-i}.
\]
Indeed its denominator lies between $dt^{d-1}$ and
$d(q+1)^{(d-1)/d}$; use $q^i-1\ge(1-1/q)q^i$
for the upper bound and $3(q-1)\ge q+1$ for the lower bound.

A selected root sum is $\sum_i a_i C_i$, where $C_i$ is a subset sum
of the $d$th roots of unity. Always $|C_i|\le d$; if $C_i\ne0$,
its nonzero cyclotomic integer norm gives
$|C_i|\ge d^{-(d-2)}$. Similarly, a nonzero $\sum_i C_i$ has absolute
value at least $(dM)^{-(d-2)}$, since every other conjugate has
absolute value at most $dM$. In that case
\[
 \left|\sum_i\delta_i C_i\right|
 \le t\sum_{i\ge L}q^{-i}<2tq^{-L}
 <\tfrac12t(dM)^{-(d-2)},
\]
so it cannot cancel $t\sum_i C_i$.
If $\sum_i C_i=0$, take the least $i$ with $C_i\ne0$.
Its perturbation has absolute value greater than
$tq^{-i}/(3d^{d-1})$, while the entire later tail has absolute value
at most $tq^{-i}/(q-1)$, strictly smaller by hypothesis.
Thus all $C_i$ must vanish. For prime $d$, irreducibility of $\Phi_d$
then forces each selected subset to be empty or a whole orbit.
A singleton cannot qualify.
\end{proof}

\begin{theorem}[Unique witnesses against a doubled exponent]
\label{tou:theorem}
Fix a rational rate $\rho=a/b\in(0,1)$ in lowest terms and a prime
$d\nmid a$. For arbitrarily large exact-rate lengths $n$, and over
infinitely many sufficiently large prime fields for each such length,
there are Reed--Solomon codes and received lines with gap
\[
 \eta=\frac{2d+1}{n}
\]
and exactly $J=2^{(H_2(\rho)/d+o(1))n}$ nearby parameters.
Every nearby point has a unique codeword at exact distance
$\theta=1-\rho-\eta$, while parameter zero has exact distance
\[
 1-\rho-1/n=\theta+\frac{2d}{2d+1}\eta.
\]
The radius is strictly below Elias and correlated agreement fails.
For every fixed $c_1>0$ and $0<c_2<2+1/d$,
\[
 \frac{J}{c_1n2^{c_2H_2(\rho)/\eta}}=2^{\Omega_{\rho,d,c_2}(n)}.
\]
Given the specified finite-field roots, the code, line, and a decoder
for every nearby parameter can be constructed in polynomial time in
$n+\log p$. No comparable bound on finding a suitable splitting prime
is asserted.
In particular, $c_2=2$ is compatible with separation at least
$(1-\varepsilon)\eta$ for any fixed $\varepsilon>0$.
\end{theorem}
\begin{proof}
Choose $n=bu$ with $au\equiv-1\pmod d$, let $c$ be its residue in
$\{0,\ldots,d-1\}$, and set
\[
 m=(n-c)/d-2,\qquad D=(\rho n+1)/d,\qquad M=m+c.
\]
For large $n$ we have $1\le D<m$. Choose $q,L$ as in
Lemma~\ref{tou:rigidity}. Use the first $m$ radical orbits as core,
the next $c$ radical points as extra coordinates, and the two padding
orbits $\mu_d$ and $t\mu_d$. These are distinct in characteristic zero,
as their magnitudes are distinct except within each orbit.

\emph{Reduce without new root relations.}
Let $E$ be the Galois field generated by $\zeta,t$ and all $a_i$.
Its degree is at most $N=(d-1)d^{M+1}$. Put
\[
 W=\prod_{i=L}^{L+M-1}(q^i-1),\qquad
 B=\bigl(Wn(q+1)\bigr)^N.
\]
All displayed coordinates become algebraic integers after multiplication
by $W$. A nonzero difference of two coordinates or a nonzero subset sum
of nonpadding coordinates has, after this multiplication, nonzero integer
norm of absolute value at most $B$: each conjugate has absolute value
at most $Wn(q+1)$. Choose a completely split prime $p>B$.
The finite-field reductions are therefore distinct and retain exactly
the root-sum rigidity just proved. Arbitrarily large such primes exist
by the splitting-prime fact used in Lemma~\ref{fg:lift}; only infinitude
is needed. We may also require
\[
 p>q^{S_{\max}},\qquad
 S_{\max}=D(2L+2m-D-1)/2.
\]

\emph{Keep one index-sum class.}
Among the $D$-subsets of $\{L,\ldots,L+m-1\}$, choose a sum $S$ of
maximum multiplicity. Its class has size
\[
 J\ge\frac{\binom mD}{D(m-D)+1}.
\]
For every support $I$ put
$H_I(X)=\prod_{i\in I}(X^d-R_i)$. Fix any reference support of size
$D$, set $w=H_{I_0}$, and let $f$ be its evaluation word. Set the
line direction to zero on the nonpadding coordinates, to $q^S$ on
$\mu_d$, and to one on $t\mu_d$.
The identities
\[
 \frac{H_I(1)}{H_I(t)}=q^{\sum_{i\in I}i},\qquad
 -H_I(t)=\frac{(-1)^{D+1}(q-1)^D}{\prod_{i\in I}(q^i-1)}
\]
show that $w-H_I$ is a nearby codeword whenever $\sum I=S$.
Its degree is at most $dD-d<K=dD-1$, and it matches $K+1$
core coordinates and all $2d$ padding coordinates.
The root bound and the support of the direction give the exact distances.

Conversely, any nearby codeword must match all padding and $K+1$
nonpadding coordinates. Its monic difference from $w$ has degree
$dD$ and zero penultimate coefficient. Root-sum rigidity forces its
roots to be $D$ whole core orbits. The two padding values then imply
$q^{\sum I}=q^S$ in $\F_p$. Since both integer powers are less than
$p$, the index sums are equal. This classifies every nearby codeword.

\emph{Recover the unique witness.}
For a nonzero parameter $z$, lift
$(-1)^{D+1}(q-1)^D/z$ to $\{0,\ldots,p-1\}$.
Every genuine nearby label yields
$T_I=\prod_{i\in I}(q^i-1)<q^{S_{\max}}<p$.
The base-$q$ greedy decoder from Lemma~\ref{pd:integer-products}
recovers the support uniquely. Check its cardinality and index sum;
rejection proves the parameter is outside the tested radius.
Dynamic programming finds a largest index-sum class in polynomial time
in $m$; the integer products, finite-field locators, and decoder use
polynomial time in $n+\log p$.

Finally $m=n/d+O_d(1)$ and $D=\rho n/d+O(1)$, so both the lower
bound on $J$ and the trivial upper bound $J\le\binom mD$ give
$\log_2J=nH_2(\rho)/d+O_{\rho,d}(\log n)$.
Comparison with $c_2nH_2(\rho)/(2d+1)$ proves the numerical ratio.
The splitting primes can be chosen arbitrarily large, ensuring strict
Elias. Choosing a large prime $d$ not dividing $a$ proves the final
separation statement. The gap tends to zero, and the nearby set is sparse.
\end{proof}
